\documentclass[11pt]{article}
\usepackage[utf8]{inputenc}
\usepackage{graphicx}
\usepackage{amsmath}
\usepackage{booktabs}
\usepackage{multirow}
\usepackage{algorithm}
\usepackage{algorithmic}
\usepackage{cite}
\usepackage[tableposition=top]{caption}

\title{Hierarchical Fault Detection and Classification for Industrial Bearing Diagnosis: A Safety-First Approach}

\author{Anonymous Authors}

\date{\today}

\begin{document}

\maketitle

\begin{abstract}
Industrial bearing fault diagnosis requires not only high accuracy but also \emph{safety guarantees} --- failing to detect a fault (false negative) is far more costly than a false alarm. This paper proposes a hierarchical classification framework that mirrors the actual maintenance decision process: Level 1 determines whether the bearing is faulty (Yes/No), while Level 2 identifies the specific fault type. We discover that entropy-based features (sample entropy, permutation entropy) achieve \textbf{100\% fault detection rate} at Level 1, making them ideal for safety-critical fault detection. However, these entropy features perform only marginally better than random guessing at Level 2 (33\% accuracy for three fault types), revealing a fundamental limitation: entropy captures signal regularity but not fault-type-specific patterns. This finding motivates an adaptive feature selection strategy where different feature subsets are optimal for different classification tasks. Experimental results on the CWRU bearing dataset demonstrate that the hierarchical approach with entropy features provides guaranteed fault detection while maintaining interpretable decision-making for maintenance prioritization.
\end{abstract}

\section{Introduction}

Rotating machinery failures in industrial settings can lead to catastrophic equipment damage, production losses, and safety hazards. Bearings are among the most critical components, and their failure accounts for a substantial portion of rotating machinery breakdowns. Early and accurate fault detection is essential for predictive maintenance, yet real-world deployments face a fundamental tension: the cost of a missed fault (false negative) vastly outweighs the cost of a false alarm.

In practice, maintenance decisions follow a hierarchical process:

\begin{enumerate}
    \item \textbf{Is the bearing faulty?} (Yes/No) --- The critical safety decision
    \item \textbf{What type of fault?} (Ball/Inner Race/Outer Race) --- The diagnostic decision for repair planning
\end{enumerate}

This two-stage hierarchy reflects how maintenance engineers actually work: first, they must ensure the system is safe, then they need to know \emph{what} to repair.

Existing approaches treat fault diagnosis as a flat multi-class classification problem, directly predicting one of $k$ fault types. However, this formulation ignores the asymmetric costs of errors: a fault missed at the first stage propagates through the entire system, while errors at the second stage merely affect repair efficiency.

\textbf{Our key observation:} Entropy features (sample entropy, permutation entropy) excel at determining whether a fault exists, but they fail to distinguish between fault types. This observation motivates our hierarchical approach:

\begin{itemize}
    \item \textbf{Level 1 (Safety)}: Use entropy features for guaranteed fault detection
    \item \textbf{Level 2 (Diagnosis)}: Requires richer features (time/frequency domain) for fault type classification
\end{itemize}

\textbf{Contributions:}

\begin{itemize}
    \item We propose a hierarchical classification framework aligned with practical maintenance workflows, explicitly separating the safety decision (fault detection) from the diagnostic decision (fault typing).

    \item We demonstrate that entropy features achieve 100\% fault detection rate on the CWRU bearing dataset, making them ideal for Level 1 safety-critical detection.

    \item We reveal a fundamental limitation of entropy features: they achieve only random-level accuracy (33\%) for three-class fault type classification, suggesting they capture signal regularity but not fault-type-specific patterns.

    \item We analyze the implications for adaptive feature selection: different classification tasks benefit from different feature subsets, motivating future work on task-aware feature switching.
\end{itemize}

\section{Related Work}

\subsection{Imbalance Learning in Fault Diagnosis}
Class imbalance is pervasive in industrial fault diagnosis. While normal operational data accumulates continuously, fault samples are relatively rare. Traditional approaches include random undersampling (RUS) of the majority class and oversampling techniques such as SMOTE. However, these methods treat all classification errors as equal, ignoring the asymmetric costs in practical maintenance scenarios.

\subsection{Hierarchical Classification}
Hierarchical classification has been studied in text categorization and medical diagnosis, where class labels form a hierarchy. However, its application to industrial fault diagnosis remains underexplored. Existing work typically treats the problem as flat multi-class classification, missing the opportunity to align the classifier with the actual decision-making process.

\subsection{Entropy Features for Fault Detection}
Entropy measures have been widely used for fault detection in rotating machinery. Sample entropy quantifies signal regularity, while permutation entropy captures the complexity of time series. These features are particularly effective for detecting the presence of anomalies but are less studied for distinguishing between fault types.

\section{Methodology}

\subsection{Problem Formulation}
Given a set of vibration signals, we formulate fault diagnosis as a two-stage hierarchical classification:

\begin{enumerate}
    \item \textbf{Level 1 (Fault Detection)}: Binary classification
    \[
    y_1 \in \{0, 1\} \quad \text{(0 = Normal, 1 = Fault)}
    \]
    The goal is to minimize missed faults (false negatives) with guaranteed detection rate.

    \item \textbf{Level 2 (Fault Typing)}: If $y_1 = 1$, predict the fault type
    \[
    y_2 \in \{1, 2, 3\} \quad \text{(1 = Ball, 2 = Inner Race, 3 = Outer Race)}
    \]
    The goal is to accurately identify the fault type for repair planning.
\end{enumerate}

This formulation explicitly models the maintenance decision process and allows task-specific feature optimization.

\subsection{Experimental Design: Simulating Industrial Imbalance}

In real industrial settings, normal operational data is abundant while fault data is scarce. To simulate this scenario while maintaining sufficient training data for Level 2 classification:

\begin{itemize}
    \item \textbf{Normal samples}: Bootstrap resampling from the available normal data to create $n_{normal} = n_{fault} \times IR$ samples, simulating the abundance of normal data in production environments.
    \item \textbf{Fault samples}: Retained at a fixed number (100 per class) to ensure adequate training data for Level 2 classification.
    \item \textbf{Imbalance Ratio (IR)}: Controlled as $IR = n_{normal} / n_{fault}$, with values 3, 5, and 10 tested.
\end{itemize}

This design reflects the practical scenario where: (1) normal data is plentiful during operation, and (2) historical fault records provide enough examples for training but are relatively scarce compared to normal data.

\subsection{Meta-Feature Extraction}
We extract three categories of features from vibration signals:

\begin{table}[htbp]
\centering
\caption{Meta-Feature Summary}
\label{tab:features}
\begin{tabular}{|l|l|l|}
\toprule
Category & Feature & Physical Meaning \\ \midrule
Time-Domain & RMS, Variance, Peak-to-peak, etc. & Signal energy and variation \\ \midrule
Frequency-Domain & Main Freq. Ratio, Spectral Entropy, etc. & Frequency distribution characteristics \\ \midrule
Entropy & Sample Entropy, Permutation Entropy & Signal regularity and complexity \\
\bottomrule
\end{tabular}
\end{table}

\textbf{Key hypothesis:} Entropy features capture the \emph{presence} of faults (abnormal regularity) but not the \emph{type} of faults (which require frequency-domain signatures).

\section{Experimental Results}

\subsection{Dataset and Setup}
Experiments are conducted on the Case Western Reserve University (CWRU) bearing dataset with 4 classes: Normal, Ball (B), Inner Race (IR), and Outer Race (OR). Bootstrap resampling is used to generate normal samples, creating datasets with IR = 3, 5, and 10 while maintaining sufficient fault samples (100 per class) for training.

\subsection{Level 1: Fault Detection Performance}

\begin{table}[htbp]
\centering
\caption{Level 1 (Fault Detection) Results}
\label{tab:level1}
\begin{tabular}{|l|c|c|c|c|}
\toprule
IR & Recall (Fault Detection) & F1-Score & False Alarm Rate & Normal Samples \\ \midrule
3 & \textbf{1.0000} & 1.0000 & 0.0000 & 300 \\
5 & \textbf{1.0000} & 1.0000 & 0.0000 & 500 \\
10 & \textbf{1.0000} & 1.0000 & 0.0000 & 1000 \\
\bottomrule
\end{tabular}
\end{table}

\textbf{Key finding:} Entropy features achieve 100\% fault detection rate across all IR configurations. This confirms that entropy effectively captures the \emph{anomalous regularity} present in fault signals, making it ideal for safety-critical Level 1 detection.

\subsection{Level 2: Fault Type Classification Performance}

\begin{table}[htbp]
\centering
\caption{Level 2 (Fault Type Classification) Results}
\label{tab:level2}
\begin{tabular}{|l|c|c|c|c|}
\toprule
IR & Accuracy & F1-macro & Fault Samples & Random Baseline \\ \midrule
3 & 0.343 & 0.339 & 300 & 0.333 \\
5 & 0.323 & 0.318 & 300 & 0.333 \\
10 & 0.343 & 0.339 & 300 & 0.333 \\
\bottomrule
\end{tabular}
\end{table}

\textbf{Key finding:} Entropy features perform only marginally better than random guessing (33\%) for fault type classification. This reveals that while entropy captures the \emph{existence} of faults, it fails to capture fault-type-specific patterns.

\subsection{Comparison with Flat Classification}

\begin{table}[htbp]
\centering
\caption{Hierarchical vs Flat Classification}
\label{tab:comparison}
\begin{tabular}{|l|c|c|c|c|}
\toprule
IR & Flat F1-macro & Hierarchical Fault Detection & Hierarchical Fault Type Acc. \\ \midrule
3 & 0.593 & \textbf{1.0000} & 0.343 \\
5 & 0.616 & \textbf{1.0000} & 0.323 \\
10 & 0.633 & \textbf{1.0000} & 0.343 \\
\bottomrule
\end{tabular}
\end{table}

While the flat classifier achieves higher overall F1-macro, it provides no guarantee on fault detection rate. In safety-critical applications, the hierarchical approach's 100\% detection rate is preferable to the flat approach's 59-63\% F1-macro.

\section{Discussion}

\subsection{Why Entropy Works for Level 1}
Entropy measures (sample entropy, permutation entropy) quantify the regularity and complexity of time series. Fault signals exhibit different regularity patterns compared to normal signals:

\begin{itemize}
    \item \textbf{Normal signals}: High regularity (periodic, predictable) $\rightarrow$ Low entropy
    \item \textbf{Fault signals}: Low regularity (chaotic, unpredictable) $\rightarrow$ High entropy
\end{itemize}

This physical interpretation explains why entropy features excel at detecting the \emph{presence} of faults.

\subsection{Why Entropy Fails for Level 2}
Distinguishing between fault types (Ball vs. Inner Race vs. Outer Race) requires understanding the \emph{frequency signatures} of different fault mechanisms:

\begin{itemize}
    \item Ball defects: High-frequency impulses at ball passage frequency
    \item Inner race defects: Characteristic frequency related to inner race rotation
    \item Outer race defects: Characteristic frequency related to outer race fixed position
\end{itemize}

These frequency-domain characteristics are not captured by entropy measures, explaining why Level 2 accuracy is near random.

\subsection{Implications for Feature Selection}
Our findings suggest an \emph{adaptive feature selection} strategy:

\begin{itemize}
    \item \textbf{Level 1 (Safety)}: Use entropy features $\rightarrow$ Guaranteed fault detection
    \item \textbf{Level 2 (Diagnosis)}: Use frequency-domain features $\rightarrow$ Accurate fault typing
\end{itemize}

This aligns with the principle that \emph{different classification tasks benefit from different feature representations}.

\section{Conclusion}

This paper proposed a hierarchical classification framework for industrial bearing fault diagnosis that mirrors the practical maintenance decision process. Our key findings are:

\begin{enumerate}
    \item \textbf{Entropy features enable 100\% fault detection}: Sample entropy and permutation entropy reliably distinguish normal from faulty signals, achieving perfect recall across all IR configurations.

    \item \textbf{Entropy features fail for fault typing}: Achieving only 33\% accuracy (random level) for three-class fault type classification, entropy features do not capture fault-type-specific frequency signatures.

    \item \textbf{Hierarchical classification provides safety guarantees}: By separating the safety decision (Level 1) from the diagnostic decision (Level 2), the framework ensures fault detection is never compromised for classification accuracy.

    \item \textbf{Adaptive feature selection is needed}: Different classification tasks require different features. Future work should explore automatic feature switching based on the classification task.
\end{enumerate}

This work contributes to the understanding of feature-task alignment in industrial fault diagnosis, providing guidance for practitioners on when to use entropy-based vs. frequency-domain features.

Future work will explore:
\begin{itemize}
    \item Combining entropy features (Level 1) with frequency-domain features (Level 2) in a unified framework
    \item Cost-sensitive learning that explicitly models the asymmetric costs of false negatives vs. false positives
    \item Online adaptation for concept drift in industrial monitoring
\end{itemize}

\section*{Reproducibility Statement}
The source code and experimental data are available at \url{https://github.com/andyhan100044/meta-feature-resampling}.

\bibliography{references}
\bibliographystyle{plain}

\end{document}
